%% kapa: exam=1 round=31 year=2020 period=1 form=A subject=civil_law subject_ko=민법 seq=28 number=28
%% source: Q-Net 감정평가사 기출문제 (cst00309) artlSeq=5207219 file=extracted/1008/5207219/A형/A형/2020년도 제31회 감정평가사 1차 시험_1교시_A형.pdf
%% conversion: kordoc 4.24.51 -> tools/gen1.py (auto); review=proofread-sonnet; flags=none
\begin{question}{28}
점유자와 회복자의 관계에 관한 설명으로 옳지 않은 것은?
\begin{choices}
  \item 선의의 점유자는 점유물의 과실을 취득한다.
  \item 과실의 수취에 관하여 점유자의 선․악의는 과실이 원물에서 분리되는 때를 기준으로 판단한다.
  \item 악의의 점유자는 그가 소비한 과실의 대가를 보상하여야 한다.
  \item 그의 책임있는 사유로 점유물을 멸실ㆍ훼손한 선의의 타주점유자는 손해 전부를 배상하여야 한다.
  \item 과실을 취득한 점유자는 그가 지출한 비용 전부를 청구할 수 있다.
\end{choices}
\end{question}
\answer{5}
